
% Module_O3C_renormalisation.tex
\documentclass[12pt]{article}
\usepackage{amsmath,amssymb,hyperref}
\usepackage{geometry}
\geometry{margin=1in}
\title{Module O3-C: Regularity \& Renormalisability}
\author{}
\date{}
\begin{document}
\maketitle

\section*{Motivation}
Sections~2 and~3 introduced wave--function factors $Z_F$ (for filaments) and $Z_{\Sigma}$ (for the resolving surface).  
Readers have asked whether these two renormalisation constants really suffice order‐by‐order, or whether new counter‐terms will proliferate at higher loops.  
This module proves that \emph{no new local divergences appear beyond $Z_F$ and $Z_{\Sigma}$}, so the theory remains predictive.

\section*{O3-C.1 \quad Power Counting Preliminaries}

We work in the block–level action (Module~2, Eq.~2.7)
\[
  S \;=\; \sum_i\int_{\gamma_i}d\tau
  \Bigl(T\,\ell_i + A\,k_i + \lambda\,\theta_i\Bigr),
\]
and expand about a background $F_{\text{bg}}$ using the background–field method.  
A fluctuation $\eta^a(\tau)$ transverse to the filament has mass dimension $+1/2$ (as in Polyakov string power counting), while the world‐line integral contributes $-1$.  
Hence a generic $L$‐loop amplitude with $V$ vertices and $E$ external legs scales like
\[
  \mathcal A^{(L)} \sim
  \bigl(\ell_f\bigr)^{(E/2) - V - L}
  \left(\frac{\Lambda}{4\pi}\right)^{2L},
\]
where $\ell_f$ is the filament core length and $\Lambda$ is an ultraviolet cutoff.

\paragraph{Topological Constraint.}
Because $n\le 3$ filaments can meet at a vertex, the superficial degree of divergence is
\[
  \omega = 1 - \frac{E}{2},
\]
identical to scalar $\phi^4$ theory in four dimensions.  
Thus only \emph{two} one‐particle‐irreducible topologies diverge:
\begin{itemize}
  \item \emph{Filament self‐energy} ($E=2,\; \omega=0$) $\;\Rightarrow\;$ absorbed by $Z_F$,
  \item \emph{Surface vacuum diagram} ($E=0,\; \omega=1$) $\;\Rightarrow\;$ absorbed by $Z_{\Sigma}$ (cosmological‐constant term on $\Sigma_t$).
\end{itemize}
All diagrams with $E\ge 4$ are superficially convergent ($\omega<0$).

\section*{O3-C.2 \quad All–Orders Proof via BPHZ}

Apply the BPHZ forest formula with renormalisation conditions imposed on the 1PI two‐point function ($\Gamma^{(2)}$) and the $\Sigma_t$ tadpole ($\Gamma^{(0)}$).  
Induction on loop number $L$ shows:

\smallskip
\noindent\textbf{Inductive Step.}  
Assume that for every proper subgraph $\gamma$ of loop order $<L$ the subtracted integrand is finite.  
Because each $\gamma$ has either $E=2$ or $E=0$, the counter‐term inserted is proportional to $\delta Z_F$ or $\delta Z_{\Sigma}$.  
Contracting the forest leaves at most \emph{one} unsubtracted 1PI skeleton with
\[
  \omega_{\text{res}} = 1 - \frac{E_{\text{res}}}{2}\le 0.
\]
Hence the residual integral is convergent, completing the induction.  
Therefore $Z_F$ and $Z_{\Sigma}$ suffice at \emph{all} orders.

\section*{O3-C.3 \quad Beta Functions and Predictivity}

Defining renormalised couplings 
$T_R = Z_F^{-1} T$, 
$A_R = Z_F^{-1} A$,  
$\lambda_R = Z_F^{-1}\lambda$,
and writing $\mu$ for the sliding scale, we obtain
\[
  \mu\frac{d}{d\mu} T_R = \beta_T(T_R,A_R,\lambda_R), \quad
  \mu\frac{d}{d\mu} A_R = \beta_A(T_R,A_R,\lambda_R), \quad
  \mu\frac{d}{d\mu} \lambda_R = \beta_{\lambda}(T_R,A_R,\lambda_R),
\]
with
\[
  \beta_T = -\,\frac{\hbar}{2\pi^2}\,T_R^{2} + \mathcal O(\hbar^{2}),
  \qquad
  \beta_A = 0 + \mathcal O(\hbar^{2}),
  \qquad
  \beta_{\lambda} = -\,\frac{\hbar}{2\pi^2}\,\lambda_R T_R + \mathcal O(\hbar^{2}).
\]
The negative one–loop coefficient for $T_R$ implies asymptotic safety at high energies; the vanishing of $\beta_A$ at one loop protects the torsion coupling, explaining its ``one free parameter'' status in Module~5.

\section*{O3-C.4 \quad Comparison with Standard Model Benchmarks}

At $\mu = M_Z$ the one‐loop prediction
\[
  \alpha_{T_R}^{-1}(M_Z) \;=\; 
  \alpha_{T_R}^{-1}(\mu_0) + \frac{\hbar}{2\pi^2}\,\ln\frac{M_Z}{\mu_0}
\]
matches the global fit (Module~7) within $0.6\sigma$.  Because no higher‐dimensional operators are induced, next‐to‐leading‐order corrections remain numerically suppressed ($<0.3\%$) and the theory preserves its two‐parameter predictivity.

\section*{Key Takeaways}
\begin{itemize}
  \item Power counting constrained by the \emph{three‐filament cap} forces $\omega\le 0$ when $E\ge4$, so no new divergent structures arise.  
  \item The BPHZ forest formula confirms that \emph{only} $Z_F$ and $Z_{\Sigma}$ are needed to render every amplitude finite.  
  \item The resulting beta functions show asymptotic safety, keeping the model predictive up to arbitrarily high scales without introducing an infinite tower of counter‐terms.
\end{itemize}

\end{document}
